\chapter{Background and Related Work}
\label{ch:background}

\section{RISC-V Instruction Set Architecture}
\label{sec:riscv-isa}

\subsection{History and Design Philosophy}

The RISC-V project began at the University of California, Berkeley in 2010 under
the direction of Krste Asanović and David Patterson, motivated by the observation
that every major instruction set architecture either carried decades of backwards-
compatibility baggage (x86, ARM) or was encumbered by licensing terms that inhibited
academic and open-source development \cite{patterson2020risc}. The design team
deliberately started from a clean slate, taking the well-understood principles of
reduced instruction set computing and applying them with the benefit of thirty years
of hindsight.

The central tenets of the RISC-V philosophy are modularity, simplicity, and openness.
Rather than defining a single monolithic instruction set, RISC-V specifies a small
mandatory base ISA (RV32I or RV64I) and a collection of standard extensions that can
be combined selectively. A chip targeting IoT sensor fusion needs only RV32I and
perhaps the C (compressed) extension; a high-performance server processor adds the M,
A, F, D, and V extensions. The letter ``G'' denotes the general-purpose combination
\texttt{IMAFD}. This modularity allows silicon area to be allocated precisely to the
features that a given application domain demands.

The open-source nature of RISC-V distinguishes it from all commercial alternatives.
The specifications are published under Creative Commons licensing, the reference
implementations (Rocket, BOOM, CVA6) are available under BSD or Apache licenses, and
a rich ecosystem of open-source tools --- GCC, LLVM, Spike ISA simulator, Verilator,
and ModelSim flows --- eliminates the toolchain cost that historically made ASIC design
prohibitive for academic groups.

\subsection{Base ISA: RV32I}

RV32I defines thirty-two 32-bit general-purpose integer registers (\texttt{x0}--\texttt{x31},
where \texttt{x0} is hardwired to zero), a 32-bit program counter, and a load/store
memory architecture in which all arithmetic operates on registers and memory is
accessed exclusively through explicit load and store instructions \cite{riscv-spec-2019}.
The instruction encoding is fixed at 32 bits (the C extension provides optional 16-bit
compressed forms), and all instructions are aligned on 4-byte boundaries in the
standard configuration.

The RV32I instruction formats --- R, I, S, B, U, and J --- are designed to keep the
source and destination register fields at fixed bit positions (bits [11:7] for \texttt{rd},
[19:15] for \texttt{rs1}, [24:20] for \texttt{rs2}) regardless of format, reducing
decoder complexity. The opcode occupies bits [6:0] and the \texttt{funct3}/\texttt{funct6}/\texttt{funct7}
fields further disambiguate within an opcode class.

The base integer instruction set covers arithmetic (\texttt{ADD}, \texttt{SUB},
\texttt{ADDI}), logical (\texttt{AND}, \texttt{OR}, \texttt{XOR} and their immediate
forms), shift (\texttt{SLL}, \texttt{SRL}, \texttt{SRA}), comparison (\texttt{SLT},
\texttt{SLTU}), branch (\texttt{BEQ}, \texttt{BNE}, \texttt{BLT}, \texttt{BGE},
\texttt{BLTU}, \texttt{BGEU}), unconditional jump (\texttt{JAL}, \texttt{JALR}),
and memory access (\texttt{LW}, \texttt{LH}, \texttt{LB}, \texttt{LHU}, \texttt{LBU},
\texttt{SW}, \texttt{SH}, \texttt{SB}).

\subsection{M Extension: Integer Multiplication and Division}

The M extension adds hardware multiply and divide instructions to RV32I
\cite{riscv-spec-2019}. \texttt{MUL} produces the lower 32 bits of the product;
\texttt{MULH}, \texttt{MULHU}, and \texttt{MULHSU} produce the upper 32 bits for
signed$\times$signed, unsigned$\times$unsigned, and signed$\times$unsigned cases
respectively. \texttt{DIV}, \texttt{DIVU}, \texttt{REM}, and \texttt{REMU} perform
integer division and remainder with signed and unsigned variants. Without the M
extension, multiplication must be emulated in software using shift-and-add sequences,
which is impractical for the matrix multiplication and convolution kernels that motivate
this design.

\subsection{Privileged Specification and CSRs}

The Zicsr extension defines a uniform mechanism for accessing \glspl{csr} using the
\texttt{CSRRW}, \texttt{CSRRS}, \texttt{CSRRC}, and their immediate variants
(\texttt{CSRRWI}, \texttt{CSRRSI}, \texttt{CSRRCI}) \cite{riscv-spec-2019}. CSR
addresses occupy a 12-bit field in the instruction encoding (\texttt{inst[31:20]}),
giving a 4096-entry address space partitioned by privilege level and read/write
access. The vector extension uses this mechanism to expose three user-level CSRs:
\texttt{vl} (current vector length, address \texttt{0xC20}), \texttt{vtype} (vector
type register encoding SEW, LMUL, and tail/mask policy, address \texttt{0xC21}), and
\texttt{vlenb} (vector register length in bytes, read-only constant, address \texttt{0xC22}).

\section{Vector Processing Fundamentals}
\label{sec:vector-fundamentals}

\subsection{Flynn's Taxonomy and Parallelism Classes}

Flynn's taxonomy \cite{harris2012digital} classifies computer architectures by the
number of instruction streams and data streams they operate on simultaneously:
Single Instruction Single Data (SISD) describes a classical scalar processor; Multiple
Instruction Multiple Data (MIMD) describes a multi-core system; Single Instruction
Multiple Data (SIMD) describes the class that encompasses both vector processors and
modern packed-SIMD extensions.

Within the \gls{simd} class, a useful distinction separates \emph{true vector} processors
from \emph{packed SIMD} designs. In packed SIMD (MMX, SSE, AVX, NEON), the vector
register width is fixed in the \gls{isa} and the programmer explicitly handles loop
tails when the data set length is not a multiple of the fixed vector width. In a true
vector processor, the \gls{isa} is parameterised by a hardware-configured vector length
that the application queries at runtime using a configuration instruction; the hardware
handles tail elements automatically. RISC-V RVV follows the true vector model.

\subsection{Key Vector Concepts}

\paragraph{Vector Length and VLEN.}
\gls{vlen} is the number of bits in each vector register. For a given VLEN and SEW
(the number of bits per element), the number of elements per register is
$\text{VLEN} / \text{SEW}$. This thesis fixes VLEN=128 bits, so a register holds
sixteen 8-bit elements, eight 16-bit elements, or four 32-bit elements when LMUL=1.

\paragraph{Standard Element Width (SEW).}
\gls{sew} specifies the data type width for a vector instruction. The RVV specification
encodes SEW in a three-bit field \texttt{vsew[2:0]} within the \texttt{vtype} CSR.
Zve32x mandates support for SEW=8, 16, and 32 bits.

\paragraph{Length Multiplier (LMUL).}
\gls{lmul} groups consecutive vector registers into a single logical register, scaling
the effective vector length by the multiplier. LMUL=1 is the base case; LMUL=2 pairs
two registers, doubling the number of elements that a single instruction operates on.
LMUL=4 and LMUL=8 continue the progression. Fractional LMUL values (LMUL=1/2, 1/4,
1/8) use only a fraction of a register and are defined in the full V extension but not
required by Zve32x. The maximum vector length for a given SEW and LMUL is:
\begin{equation}
  \text{VLMAX} = \left\lfloor \frac{\text{LMUL} \times \text{VLEN}}{\text{SEW}} \right\rfloor
  \label{eq:vlmax}
\end{equation}

\paragraph{Masking.}
RVV provides a per-element masking mechanism. When the \texttt{vm} bit in an instruction
is 0, each element operation is conditioned on the corresponding bit of vector register
\texttt{v0}. Elements whose mask bit is 0 are either left unchanged (mask-undisturbed
policy) or written with all-ones (mask-agnostic policy), depending on the \texttt{vma}
bit in \texttt{vtype}. Masking enables vectorised control flow: a conditional operation
on a subset of elements can be expressed as a comparison followed by a masked arithmetic
instruction without branching.

\paragraph{Tail Behaviour.}
When vl is set to a value less than VLMAX, the elements from index vl to VLMAX-1
are ``tail'' elements. The tail-agnostic (\texttt{ta}) bit in vtype controls whether
tail elements are left unchanged or overwritten with all-ones. This thesis implements
tail-agnostic behaviour, writing zero to tail positions in the register file.

\subsection{Performance Advantages of Vector Processing}

The performance advantage of vector processing over scalar processing for data-parallel
workloads stems from three sources. First, instruction-fetch overhead is amortised:
a single \texttt{vadd.vv} instruction replaces $N$ scalar \texttt{add} instructions,
reducing instruction cache pressure proportionally. Second, branch prediction overhead
is eliminated because the innermost loop over elements is implicit in the vector length
counter inside the hardware. Third, the hardware can pipeline element operations in
a deeper pipeline than a scalar functional unit, because the latency can be hidden
behind the processing of subsequent elements within the same vector instruction.

For a concrete estimate, a BT.601 grayscale conversion applied to a 128$\times$128
image requires 16384 multiply-accumulate operations per colour channel. A scalar RV32I
core at 50~MHz completes this in approximately 98304 cycles (6 instructions per pixel).
A vector core with VLEN=128 and SEW=8 processes 16 elements per instruction; the same
workload requires only 1024 \texttt{vmulhu.vx} + \texttt{vredsum.vs} passes, a
$16\times$ reduction in instruction count assuming the bottleneck is arithmetic
throughput.

\section{RISC-V Vector Extension (RVV)}
\label{sec:rvv}

\subsection{Overview and Design Goals}

The RISC-V Vector Extension, version 1.0, was ratified by RISC-V International in
2021 \cite{rvv-spec-1-0}. Its primary design goal is to provide a single binary-compatible
vector programming model across a wide range of hardware implementations, from
microcontrollers with VLEN=64 to server processors with VLEN=4096. The key architectural
mechanism enabling this is the \texttt{vsetvl}/\texttt{vsetvli}/\texttt{vsetivli}
instruction family, which configures the hardware at runtime and returns the actual
vector length that the hardware can process in one instruction, allowing software to
loop without prior knowledge of the hardware's capabilities.

\subsection{Vector CSRs}

Three CSRs define the current vector execution context:

\begin{itemize}
  \item \textbf{vtype} (address \texttt{0xC21}): encodes \texttt{vill} (illegal vtype
        flag), \texttt{vma} (mask-agnostic), \texttt{vta} (tail-agnostic),
        \texttt{vsew[2:0]}, and \texttt{vlmul[2:0]}.
  \item \textbf{vl} (address \texttt{0xC20}): holds the current vector length as set
        by the most recent \texttt{vsetvl*} instruction.
  \item \textbf{vlenb} (address \texttt{0xC22}): a read-only constant equal to
        $\text{VLEN}/8$, i.e., the number of bytes in one vector register.
        For VLEN=128, \texttt{vlenb}=16.
\end{itemize}

\subsection{Vector Register File}

RVV defines 32 vector registers, \texttt{v0}--\texttt{v31}, each VLEN bits wide.
Register \texttt{v0} has the special role of the mask register when \texttt{vm=0}.
With LMUL>1, groups of registers are treated as a single logical register; for example,
LMUL=4 means that \texttt{v0}, \texttt{v4}, \texttt{v8}, $\ldots$ are the accessible
register group addresses, and each group contains four physical VLEN-bit registers
holding $4 \times \text{VLMAX}$ elements.

\subsection{Instruction Categories}

RVV instructions are categorised by their operand encoding:

\begin{itemize}
  \item \textbf{OPIVV} (\texttt{funct3=000}): vector--vector operations (both operands
        from the vector register file).
  \item \textbf{OPIVX} (\texttt{funct3=100}): vector--scalar operations (one operand
        from a scalar integer register).
  \item \textbf{OPIVI} (\texttt{funct3=011}): vector--immediate operations (one operand
        is a 5-bit sign-extended immediate).
  \item \textbf{OPMVV} (\texttt{funct3=010}): multiply-add and integer reduction
        vector--vector operations.
  \item \textbf{OPMVX} (\texttt{funct3=110}): multiply-add vector--scalar operations.
  \item \textbf{OPFVV/OPFVF}: floating-point categories (not implemented in Zve32x).
  \item \textbf{Load/Store}: unit-stride, strided, and indexed variants (Zve32x
        requires only unit-stride).
\end{itemize}

\subsection{Zve32x Subset}

The Zve32x profile is the minimum embedded vector extension targeting integer-only
workloads at 32-bit maximum element width \cite{rvv-spec-1-0}. It requires:
VLEN$\geq$32, SEW support for 8, 16, and 32 bits, all integer arithmetic and logical
operations with VV/VX/VI encodings, integer comparison and mask operations, reduction
operations (\texttt{vredsum}, \texttt{vredmax}, \texttt{vredmin} and unsigned variants),
and unit-stride vector load/store (\texttt{vle8.v}, \texttt{vle16.v}, \texttt{vle32.v},
and matching \texttt{vse*}). It explicitly does not require floating-point operations,
strided or indexed memory accesses, or VLEN exceeding 32 bits. This makes it the
natural profile for a first-generation tapeout targeting embedded workloads.

\section{ASIC Design Methodology}
\label{sec:asic-methodology}

\subsection{RTL Design Phase}

The \gls{rtl} design phase produces synthesisable SystemVerilog (or VHDL) that
describes the logical behaviour of the circuit at the level of register transfers and
combinational functions \cite{harris2012digital, weste2010cmos}. Critical rules that
distinguish synthesisable from simulation-only RTL include:

\begin{itemize}
  \item \textbf{No \texttt{initial} blocks.} Initial blocks are interpreted by
        simulators to set power-on state but are not supported by synthesis tools.
        All registers must be initialised through reset logic.
  \item \textbf{Synchronous, active-low reset.} A single active-low \texttt{rst\_n}
        signal, sampled on the rising clock edge, is the industry standard for ASIC
        reset. Asynchronous resets complicate static timing analysis because they
        can produce glitches during assertion and deassertion.
  \item \textbf{No latches.} An \texttt{always\_comb} block that does not assign
        all outputs on all code paths infers a level-sensitive latch. Latches cause
        timing violations and are eliminated by ensuring every \texttt{always\_comb}
        block has a complete set of default assignments before any conditional logic.
  \item \textbf{No simulation tasks in synthesis paths.} \texttt{\$display},
        \texttt{\$time}, \texttt{\$monitor}, and similar system tasks are
        simulation-only; they must be confined to testbenches.
  \item \textbf{Parameterised widths.} Hardcoding bit widths to specific VLEN or
        SEW values makes re-targeting to a different configuration require manual
        search-and-replace operations that are error-prone. Using \texttt{parameter}
        and \texttt{localparam} with derived widths ensures that changing one constant
        propagates correctly.
\end{itemize}

\subsection{Logic Synthesis}

Synthesis translates the RTL description into a gate-level netlist using a target
technology library (standard cell library for ASIC or LUT library for FPGA). The
synthesiser performs several optimisation passes: technology mapping selects cells from
the library; retiming may move registers across combinational logic to balance stage
depths; and area/power optimisation reduces the gate count while respecting timing
constraints \cite{weste2010cmos}. The output is a netlist of interconnected standard
cells and a delay annotated timing model.

\subsection{Static Timing Analysis}

\gls{sta} verifies that the design meets its clock frequency specification without
running gate-level simulation. It computes the worst-case propagation delay along
every combinational path from one register's output to the next register's input,
accounting for process, voltage, and temperature (PVT) corners. A path that arrives
at a flip-flop's data input too close to the rising clock edge (setup violation) or
too quickly after a previous clock edge (hold violation) will cause a functional failure
in silicon. For the designs in this thesis, the FPGA timing analyser (Quartus Prime
TimeQuest) plays the role of STA, and the Fmax reported is the reciprocal of the
critical-path delay under slow-corner conditions.

\subsection{Place and Route}

\gls{pnr} takes the synthesised netlist and assigns each cell a physical location on
the chip (placement) and then connects them using metal routing (routing). In a full
ASIC flow, this step is performed using tools such as Cadence Innovus or Synopsys
ICC2. The routed design produces a GDS-II file, the industry standard for mask
fabrication. For FPGA, Quartus Prime Fitter performs placement onto the array of
adaptive logic modules and routes through the programmable interconnect.

\subsection{Physical Verification: DRC and LVS}

\gls{drc} checks that all geometries in the GDS comply with the foundry's manufacturing
rules: minimum feature sizes, spacing between metal layers, via densities, and antenna
rules. \gls{lvs} compares the extracted netlist from the GDS (derived by recognising
transistor structures from the layout polygons) against the schematic netlist from
synthesis, verifying that the two are topologically equivalent. Both checks are
mandatory before any design is submitted for fabrication.

\section{Related Work}
\label{sec:related-work}

\subsection{Ara: Scalable RVV Vector Processor in 22nm FD-SOI}

Ara is a scalable, application-class RVV vector processor developed at ETH Zürich and
tapeout in TSMC 22nm FD-SOI technology \cite{ara2020}. It is coupled to the CVA6
(formerly Ariane) 64-bit out-of-order RISC-V scalar core and supports the full RISC-V
V extension including floating-point. The design is parameterised by the number of
lanes, with published results for two-lane, four-lane, and eight-lane configurations.
At four lanes, Ara achieves 1~GHz+ operating frequency and delivers 12.5~GFLOP/s on
double-precision floating-point reductions.

Ara's lane architecture assigns a fixed set of vector registers to each lane, enabling
all lanes to process their assigned element slice in parallel. The register file in each
lane stores VLEN/NLANE bits. An operand requester unit manages read-after-write hazards
between vector instructions through a scoreboard. The memory system uses a banked
scratchpad SRAM with one bank per lane to avoid memory bottlenecks.

Ara represents the current state of the art for high-performance RVV implementations
but is significantly more complex than a first-generation tapeout target: its lane
count parameterisation, out-of-order scalar core coupling, floating-point support, and
full V extension compliance require a design team of several engineers and a budget
commensurate with industrial silicon development.

\subsection{CVA6 / Ariane}

CVA6 is the scalar core paired with Ara \cite{cva6}. It is a six-stage, in-order,
single-issue 64-bit RISC-V core targeting application-class embedded Linux workloads.
Its contributions most relevant to this thesis are the demonstration that a small
team can produce a tapeout-quality RISC-V implementation in SystemVerilog with thorough
verification, and its published FD-SOI results showing 1.7~GHz at 1.1~V. The CVA6
core also serves as the reference for the scalar--vector dispatch interface that RVV
processors use: the scalar core stalls when the vector unit is busy, forwarding the
vector instruction and associated scalar register values directly to the coprocessor.
This thesis adopts the same dispatch paradigm for the scalar--VPU interface.

\subsection{Vitruvius+}

Vitruvius+ is an area-efficient RVV vector coprocessor for high-performance computing
developed at the Barcelona Supercomputing Centre \cite{vitruvius2021}. It emphasises
decoupled memory access by providing a large number of outstanding memory requests,
hiding DRAM latency for streaming kernels. Vitruvius+ integrates with RISC-V scalar
cores via the RVXIF (RISC-V External Interface) extension, an emerging standard for
coprocessor coupling that provides a more structured handshake than ad-hoc signal
naming. Its vector register file is distributed across multiple SRAM macros with
careful power gating to reduce idle-state power.

\subsection{Hwacha}

Hwacha is a decoupled vector accelerator developed at UC Berkeley as a predecessor to
RVV, using a custom ISA rather than the standardised RVV encoding \cite{hwacha2015}.
It introduced several concepts that influenced the RVV design, including the separation
of vector configuration (the \texttt{vsetvl} idiom) from execution, the decoupled
memory access architecture (where vector loads issue independently of computation),
and the notion of a vector fetch unit that sequences element-wise operations from a
single instruction. Hwacha was fabricated in TSMC 28nm and validated silicon results
showed correct execution of linear algebra kernels.

\subsection{Comparative Analysis}

Table~\ref{tab:related-work} summarises the key characteristics of the related designs
compared to this work.

\begin{table}[H]
  \centering
  \caption{Comparison of RISC-V vector processor implementations}
  \label{tab:related-work}
  \begin{tabular}{@{}lccccc@{}}
    \toprule
    \textbf{Design} & \textbf{ISA} & \textbf{VLEN} & \textbf{Lanes} & \textbf{Tech} & \textbf{Target} \\
    \midrule
    Ara \cite{ara2020}       & RVV 1.0 (full) & 512--4096 & 2--8 & 22nm FD-SOI & HPC/Server \\
    CVA6 \cite{cva6}         & RV64IMAFDC     & N/A        & N/A  & 22nm FD-SOI & App-class \\
    Vitruvius+ \cite{vitruvius2021} & RVV 1.0 & 512 & 8   & 7nm         & HPC \\
    Hwacha \cite{hwacha2015} & Custom (pre-RVV)& 2048      & 8    & 28nm        & HPC \\
    \textbf{This work}       & \textbf{Zve32x} & \textbf{128} & \textbf{1} & \textbf{FPGA (Cyclone~V)} & \textbf{Embedded} \\
    \bottomrule
  \end{tabular}
\end{table}

The fundamental distinction between this work and the related designs is the design
target. All of Ara, Vitruvius+, and Hwacha target high-performance computing workloads
in advanced technology nodes (7--28nm). They prioritise throughput, utilising multiple
parallel lanes, large VLENs, and sophisticated out-of-order scalar cores. This thesis
targets the opposite point in the design space: a minimum-area, tapeout-viable
implementation of the Zve32x subset intended for embedded edge inference. The single-lane
architecture, 128-bit VLEN, and 5-stage in-order scalar core reflect deliberate area
minimisation for a first-generation student tapeout.

The closest academic parallel is the open-source work on small embedded RVV processors,
but to the authors' knowledge, no published work explicitly documents a complete
single-lane Zve32x implementation with end-to-end functional verification on image
processing benchmarks in publicly available SystemVerilog RTL, making this a distinct
contribution to the open-source RISC-V hardware ecosystem.

\section{Summary}
\label{sec:background-summary}

This chapter has established the technical foundations required to understand the
design decisions in subsequent chapters. The RISC-V ISA provides a modular, open
framework; its RV32IM base gives the scalar core its programming model, and the
Zicsr extension provides the mechanism through which the vector unit exposes its
configuration state. Vector processing fundamentally improves data-parallel performance
by amortising instruction-fetch overhead across many elements, and the RVV Zve32x
subset defines a minimum yet complete integer vector profile suitable for embedded
inference. The ASIC design methodology sets strict synthesisability rules that are
enforced throughout the RTL. Finally, the related work survey shows that while
high-performance many-lane implementations exist, the single-lane embedded Zve32x
target of this thesis occupies a distinct and underserved point in the design space.
